% Ad Quintum Fratrem 2.5 (ad-quintum-fratrem-02-05) --- English translation
% Translated by Claude (Anthropic) from the Latin (Perseus).
% Style guide: STYLE.md.

\ciceroLetterOpener{MARCVS QVINTO FRATRI SALVTEM.}

\ciceroSection{1} I had sent you a letter earlier in which it was written that our Tullia had been betrothed to Crassipes on the day before the Nones of April [4 April]; and I had written out the rest, both public and private. Since then these things have been done. On the Nones of April [5 April], by decree of the Senate, money was decreed to Pompey for the corn-supply up to forty million sesterces; but on the same day vehement debate was held about the Campanian land, with the Senate's shouting almost as in a contio. The shortage of money and the dearness of corn made the case the more bitter. I shall not pass over even this: the Capitolines and the Mercuriales threw M. Furius Flaccus, a Roman knight and a worthless man, out of their college, with him present and lying at the feet of every one of them. As I was about to leave on the eighth before the Ides of April [6 April], I gave the betrothal-feast to Crassipes. From which dinner the excellent boy, your Quintus and mine, was missing, having been a little out of sorts. On the seventh before the Ides [7 April] I came to Quintus and saw him plainly recovered, and he had much and very humane conversation with me about the disagreements of our women. What more? --- nothing more amusing. Pomponia, however, even complained of you; but these things we shall handle face to face. When I had left the boy, I came to your building-site.

\ciceroSection{3} The work was being done by many builders. I urged on the contractor Longilius. He gave me his word that he wished to please us. The house will be excellent; for now more could be made out than we used to judge from the plan. Our own [house] too was being built quickly. On that day I dined at Crassipes's; after dinner I was carried in a litter to the gardens to Pompey. I had not been able to meet him in the morning, as he had been away. I wanted to see him, however, because I was to leave Rome the next day and because he had a journey into Sardinia. I met the man and asked him to send you back to us as soon as possible. He said at once. He was about to go, as he said, on the third before the Ides of April [11 April], to embark either at Labro or at Pisa. You, my brother, the moment he has come, do not let pass the first sailing, provided that the weather be suitable.

\ciceroSection{4} On the sixth before the Ides of April [8 April], before light, I wrote this letter; and I was on the road, to lodge that day at T. Titius's house in the Anagnine country, and to be the day after at Laterium, and from there, when I had been at Arpinum five days, to go on to the Pompeian estate, then on the way back to look in at the Cuman estate, so that, since the Nones of May has been the day proclaimed for Milo, I might be at Rome the day before the Nones, and might see you, my dearest and sweetest brother, by that day, as I hoped. The building at Arcanum, I thought, should be held off until your arrival. See, my brother, that you keep well, and come as soon as possible.
